% Source: https://artofproblemsolving.com/community/c3h1430077_inscribed_circle
% Stable identity: aops-topic-1430077-q01
Let the right angle be at $C$, the incircle have centre $I$ and touch the hypotenuse $AB$ at $T$. Tangent segments from a point to a circle are equal. From $C$ both tangent segments equal $r$ (with the two radii they form a square), and if $AT=x$, $BT=y$, then the legs are
\[
CA=r+x,\qquad CB=r+y,\qquad AB=x+y .
\]

\begin{asy}
usepackage("amsmath");
size(6cm);
defaultpen(fontsize(10pt));
pair ft(pair P,pair U,pair V){return U+dot(P-U,V-U)/dot(V-U,V-U)*(V-U);}
path ra(pair P,pair u,pair v,real s=0.2){return P+s*unit(u)--P+s*unit(u)+s*unit(v)--P+s*unit(v);}
pair C=(0,0),A=(15,0),B=(0,8),I=(3,3);
pair T=ft(I,A,B);
draw(A--B--C--cycle,linewidth(0.9));
draw(circle(I,3),blue);
draw(ra(C,A-C,B-C,0.9));
draw(I--(3,0),red+dashed); draw(I--(0,3),red+dashed); draw(I--T,red+dashed);
dot(I);dot((3,0));dot((0,3));dot(T);
label("$A$",A,SE); label("$B$",B,N); label("$C$",C,SW); label("$I$",I,dir(-45)); label("$T$",T,NE);
label("$r$",(1.5,0),S); label("$r$",(0,1.5),W);
label("$x$",(9,0),S); label("$y$",(0,5.5),W);
label("$x$",(T+A)/2,NE); label("$y$",(T+B)/2,NE);
label("\mbox{}",(min(currentpicture,user=true)+max(currentpicture,user=true))/2,black); // keeps the text after the figure black when LaTeX typesets the labels inline
\end{asy}

\textbf{(a)} The perimeter is $(r+x)+(r+y)+(x+y)=2r+2(x+y)=2\cdot4+2\cdot24=\boxed{56\text{ cm}}$. (Such a triangle exists: the legs $a,b$ satisfy $a+b=32$ and $a^2+b^2=576$, so $ab=224$ and the legs are $16\pm4\sqrt2$, both greater than $r$.)

\textbf{(b)} Here $x=5$, $y=12$, so by Pythagoras
\[
(r+5)^2+(r+12)^2=17^2\ \Longrightarrow\ 2r^2+34r-120=0\ \Longrightarrow\ r^2+17r-60=0=(r+20)(r-3),
\]
so $r=3$, the legs are $8$ and $15$, and the perimeter is $8+15+17=\boxed{40}$.
